\documentclass[11pt,a4paper]{article}
\usepackage{notes}

\title{Geophysics 3A --- Week 9 Workshop}
\author{Geothermics: Transient Heat Flow}
\date{\today}

\begin{document}
\maketitle

\section*{Workshop Objectives}
Students will:
\begin{itemize}
  \item Understand how thermal diffusivity controls timescales.
  \item Compare thermal and chemical diffusion processes.
  \item Interpret half-space and plate cooling models.
  \item Use Fourier series to conceptualize thermal evolution.
\end{itemize}

\section*{0:00--0:20 \quad Diffusivity and Thermal Timescales}
\begin{itemize}
  \item Review $\kappa = k / (\rho c_p)$.
  \item Emphasize modest variation in crustal rocks.
  \item Contrast sediments, rock, and ice.
\end{itemize}

\section*{0:20--0:45 \quad Thermal Memory in Earth Materials}
\begin{itemize}
  \item Penetration depth scaling: $z \sim \sqrt{\kappa t}$.
  \item Borehole paleothermometry vs ice-core records.
  \item Depth--resolution tradeoffs.
\end{itemize}

\section*{0:45--1:15 \quad Thermal vs Chemical Diffusion}
\begin{itemize}
  \item Both obey $\sqrt{Dt}$ scaling.
  \item Chemical diffusion often multi-component.
  \item Zoning persistence vs thermal smoothing.
\end{itemize}

\section*{1:15--1:25 \quad Break}

\section*{1:25--2:00 \quad Half-Space Cooling}
\begin{itemize}
  \item Assumptions and boundary conditions.
  \item Growth of thermal boundary layer.
  \item Oceanic lithosphere as a thermal clock.
\end{itemize}

\section*{2:00--2:30 \quad Plate Cooling and Fourier Series}
\begin{itemize}
  \item Fixed plate thickness as boundary condition.
  \item Fourier series as superposition of thermal modes.
  \item Mode decay times and geological interpretation.
\end{itemize}

\section*{2:30--3:00 \quad Stefan Problems in Geology}
\begin{itemize}
  \item Latent heat as effective heat capacity.
  \item Cooling times of dikes vs sills.
  \item When analytic solutions break down.
\end{itemize}

\end{document}